% !TEX root = ../../main.tex
\subsection{用户注册与防机器注册机制}

\subsubsection*{注册流程概述}

FoodShield 的用户注册流程在普通用户名注册基础上增加了口令认证和数学验证码两道机制，以降低机器批量注册和匿名身份滥用风险。完整注册流程如下：

% !TEX root = ../../../main.tex
% 用户注册流程图（含两处拒绝分支）
\begin{figure}[H]
\centering
\begin{tikzpicture}[
  node distance = 0.55cm and 2.2cm,
  proc/.style     = {rectangle, draw, rounded corners=3pt, minimum width=3.0cm, minimum height=0.6cm,
                     font=\small, fill=blue!8, draw=blue!50!black},
  decision/.style = {diamond, draw, aspect=2.0, minimum width=3.0cm, minimum height=0.75cm,
                     font=\small, align=center, fill=orange!15, draw=orange!60!black, inner sep=2pt},
  term/.style     = {rectangle, draw, rounded corners=8pt, minimum width=3.0cm, minimum height=0.6cm,
                     font=\small, fill=green!14, draw=green!50!black},
  reject/.style   = {rectangle, draw, rounded corners=8pt, minimum width=2.2cm, minimum height=0.58cm,
                     font=\small, fill=red!12, draw=red!50!black},
  arr/.style      = {-Stealth, thick},
  lbl/.style      = {font=\small},
]

\node[term]     (start)   {用户请求注册页面};
\node[proc,     below=0.65cm of start]   (captcha) {填写用户名、口令、数学验证码};
\node[decision, below=0.75cm of captcha] (d_cap)   {验证码\\正确？};
\node[reject,   right=2.2cm of d_cap]  (r_cap)   {返回错误，重新输入};
\node[decision, below=0.85cm of d_cap]  (d_dup)   {用户名\\已存在？};
\node[reject,   right=2.2cm of d_dup]  (r_dup)   {提示重复，重新输入};
\node[proc,     below=0.85cm of d_dup]  (hash)    {$\mathrm{SM3}(s\,\|\,\mathit{pwd})$ 存储口令哈希};
\node[proc,     below=0.55cm of hash]   (pid)     {生成 PID = $\mathrm{HMAC}$-$\mathrm{SM3}(K_{\mathrm{m}}, \mathit{uid}\,\|\,r)$};
\node[proc,     below=0.55cm of pid]    (db)      {写入 users 表（username, pid, pid\_r, hash, salt）};
\node[term,     below=0.55cm of db]     (done)    {注册成功，返回用户端};

\draw[arr] (start)   -- (captcha);
\draw[arr] (captcha) -- (d_cap);
\draw[arr] (d_cap)   -- node[lbl, left]{是} (d_dup);
\draw[arr] (d_cap)   -- node[lbl, above]{否} (r_cap);
\draw[arr] (d_dup)   -- node[lbl, left]{否} (hash);
\draw[arr] (d_dup)   -- node[lbl, above]{是} (r_dup);
\draw[arr] (hash) -- (pid);
\draw[arr] (pid)  -- (db);
\draw[arr] (db)   -- (done);

% 返回箭头：绕过 decision 节点（从 reject 右侧绕行回 captcha）
\draw[arr] (r_cap.north) -- ++(0,0.35) -| ([xshift=0.15cm]captcha.east);
\draw[arr] (r_dup.north) -- ++(0,0.4) -| ([xshift=-0.1cm]captcha.east);

\end{tikzpicture}
\caption{用户注册流程图}
\label{fig:register_flow}
\end{figure}


\begin{enumerate}[label=\arabic*., leftmargin=2em]
  \item 用户访问注册页面，前端向服务器请求验证码题目（\texttt{GET /captcha}）；
  \item 服务器随机生成一道整数四则运算题，将答案存入服务端 session，向前端返回题目字符串；
  \item 用户填写用户名、口令（两次确认）和验证码答案，提交注册请求（\texttt{POST /register}）；
  \item 服务器依次执行：验证码比对、用户名唯一性校验、口令长度校验、口令一致性校验；
  \item 校验通过后，生成随机盐值$s$，计算口令存储哈希：
    \[
      h_{\mathrm{pwd}} = \operatorname{SM3}(s \| \mathit{password})
    \]
    将$\langle\mathit{username}, h_{\mathrm{pwd}}, s\rangle$写入 \texttt{users} 表；
  \item 服务器以数据库内部 \texttt{user\_id} 为输入，调用 PID 生成模块，生成用户匿名身份 PID，写入 \texttt{users.pid} 和 \texttt{users.pid\_r}；
  \item 注册成功后自动完成登录，服务端 session 记录 \texttt{user\_id} 和 \texttt{username}，前端跳转至用户主页。
\end{enumerate}

\subsubsection*{口令存储机制}

为保护用户口令在数据库泄露场景下的安全性，系统采用加盐哈希存储方案，而非存储明文口令。具体地，每位用户注册时，服务器使用 \texttt{secrets.token\_hex(16)} 生成一个 128 位随机盐值$s$，对拼接后的字符串执行 SM3 哈希：
\[
  h_{\mathrm{pwd}} = \operatorname{SM3}(s \| \mathit{password})
\]
数据库中仅保存$\langle s,\, h_{\mathrm{pwd}}\rangle$，原始口令在计算完成后即丢弃。登录时，服务器从数据库读取$s$，对用户提交的口令重新计算哈希，并使用恒定时间比较函数与$h_{\mathrm{pwd}}$对比，避免时序侧信道攻击。

\subsubsection*{数学验证码机制}

FoodShield 采用服务端 session 管控的数学验证码机制抵御机器批量注册攻击：

\begin{enumerate}[label=\arabic*., leftmargin=2em]
  \item 每次获取验证码时，服务器从$[1, 20]$范围内随机抽取两个整数，随机选择加法或减法，生成题目字符串（如"$13 + 7 = ?$"）；
  \item 正确答案存入服务端 session（\texttt{session["captcha\_answer"]}），不在前端任何位置泄露；
  \item 注册提交时，服务器从 session 读取答案并与用户提交值比较；比对后无论成功与否，立即从 session 中删除该答案（一次性使用），防止枚举重试；
  \item 若 session 中不存在答案（如已过期或被提前消耗），注册请求直接拒绝。
\end{enumerate}

该机制能够有效阻止无交互式脚本批量注册，因为攻击者需要每次实时解析题目才能通过验证，而服务端 session 的一次性特性也排除了重放答案的可能。

\subsubsection*{注册完成后的 PID 初始化}

用户注册成功后，服务器以新分配的 \texttt{user\_id} 为基础输入，调用 PID 生成函数，将生成的 PID 及随机数$r$分别更新至 \texttt{users.pid} 和 \texttt{users.pid\_r}。从此，用户在系统中的通信身份以 PID 作为唯一标识，\texttt{username} 仅用于登录认证，不参与订单通信流程。PID 生成的详细设计见 \ref{subsec:pid}~节。
